\documentclass[aspectratio=169,11pt]{beamer}
% Shared Beamer style for practice contest solution walkthroughs.
\usepackage[utf8]{inputenc}
\usepackage[T1]{fontenc}
\usepackage{lmodern}
\usepackage{amsmath,amssymb}
\usepackage{xcolor}
\usepackage{hyperref}
\usepackage{listings}

\definecolor{CPNavy}{HTML}{172033}
\definecolor{CPBlue}{HTML}{2563EB}
\definecolor{CPGreen}{HTML}{16A34A}
\definecolor{CPOrange}{HTML}{F97316}
\definecolor{CPGray}{HTML}{64748B}
\definecolor{CPLight}{HTML}{F8FAFC}
\definecolor{CPCodeBg}{HTML}{F1F5F9}
\definecolor{CPCodeRule}{HTML}{CBD5E1}

\usetheme{default}
\usefonttheme{professionalfonts}
\setbeamertemplate{navigation symbols}{}
\setbeamersize{text margin left=0.65cm,text margin right=0.65cm}
\setbeamertemplate{itemize item}{\color{CPBlue}$\blacktriangleright$}
\setbeamertemplate{itemize subitem}{\color{CPGray}$\triangleright$}

\setbeamercolor{normal text}{fg=CPNavy,bg=white}
\setbeamercolor{frametitle}{fg=white,bg=CPNavy}
\setbeamercolor{title}{fg=white,bg=CPNavy}
\setbeamercolor{subtitle}{fg=CPLight,bg=CPNavy}
\hypersetup{colorlinks=true,urlcolor=CPBlue}

\setbeamertemplate{frametitle}{%
  \nointerlineskip%
  \begin{beamercolorbox}[wd=\paperwidth,ht=0.88cm,dp=0.22cm,leftskip=0.55cm,rightskip=0.35cm]{frametitle}
    \usebeamerfont{frametitle}\insertframetitle\hfill\small\insertframenumber/\inserttotalframenumber
  \end{beamercolorbox}%
}

\setbeamertemplate{footline}{%
  \leavevmode%
  \hbox{%
    \begin{beamercolorbox}[wd=.58\paperwidth,ht=2.4ex,dp=1ex,leftskip=.5cm]{author in head/foot}%
      \color{CPGray}\insertshorttitle
    \end{beamercolorbox}%
    \begin{beamercolorbox}[wd=.42\paperwidth,ht=2.4ex,dp=1ex,rightskip=.5cm plus1fil]{date in head/foot}%
      \hfill\color{CPGray}\insertshortauthor
    \end{beamercolorbox}%
  }%
}

\setbeamertemplate{title page}{%
  \vspace*{1.25cm}
  \begin{beamercolorbox}[wd=\paperwidth,sep=0.65cm,leftskip=0.15cm,rightskip=0.15cm]{title}
    {\color{white}\Huge\bfseries\inserttitle\par}
    \vspace{0.35cm}
    {\color{CPLight}\Large\insertsubtitle\par}
    \vspace{0.6cm}
    {\color{CPLight}\large\insertauthor\par}
  \end{beamercolorbox}
  \vfill
}

\newcommand{\sectiontag}[1]{\textbf{\color{CPBlue}#1}}
\newenvironment{tightitemize}{\begin{itemize}\small\setlength{\itemsep}{0.18cm}}{\end{itemize}}
\lstdefinestyle{cpcode}{
  language=Python,
  basicstyle=\ttfamily\tiny,
  keywordstyle=\color{CPBlue}\bfseries,
  commentstyle=\color{CPGray}\itshape,
  stringstyle=\color{CPGreen},
  backgroundcolor=\color{CPCodeBg},
  frame=single,
  framerule=0.25pt,
  rulecolor=\color{CPCodeRule},
  showstringspaces=false,
  columns=fullflexible,
  keepspaces=true,
  breaklines=true,
  breakatwhitespace=false,
  tabsize=4,
  xleftmargin=0.08cm,
  xrightmargin=0.08cm,
  aboveskip=0.08cm,
  belowskip=0.08cm
}


\title[lostmap]{Lost Map}
\subtitle{Day 4 solution walkthrough}
\author[Solution Deck]{Competitive Programming Bootcamp}
\date{}

\begin{document}

\begin{frame}[plain]
  \titlepage
\end{frame}

% BEGIN GENERATED STATEMENT INTERPRETATION
\begin{frame}{Interpret the Word Problem}
\small
\sectiontag{Statement excerpts:}
\begin{tightitemize}
  \item \textit{``reconstruct the map given only the table''}
  \item \textit{``output \$n-1\$ lines with two integers''}
\end{tightitemize}

\vspace{0.18cm}
\sectiontag{Programming interpretation:}
\begin{tightitemize}
  \item The distance table describes all-pairs shortest-path distances from an original tree.
  \item The original roads are recovered by selecting a minimum spanning tree.
  \item Kruskal's algorithm is a direct fit for building that tree from weighted edges.
\end{tightitemize}
\end{frame}
% END GENERATED STATEMENT INTERPRETATION







\begin{frame}{Problem Model}
\small
\sectiontag{Textbook connection:} CP4 4.3.2 Kruskal's Algorithm

\vspace{0.25cm}
\sectiontag{Core technique:} Minimum spanning tree with union-find

\vspace{0.25cm}
\begin{tightitemize}
  \item Villages are vertices.
  \item The distance matrix gives the weight of every possible road.
  \item We need roads that connect all villages with minimum total cost.
\end{tightitemize}
\end{frame}

\begin{frame}{How to Recognize the Approach}
\begin{tightitemize}
  \item Connecting all vertices at minimum weight is an MST problem.
  \item The input is dense, so converting the matrix to an edge list is straightforward.
  \item Kruskal's algorithm matches the edge-list representation.
\end{tightitemize}
\end{frame}

\begin{frame}{Algorithm}
\begin{tightitemize}
  \item Read the distance matrix.
  \item Create an edge list of (weight, u, v) pairs.
  \item Sort edges by increasing weight.
  \item Use union-find to add an edge only when it connects two different components.
  \item Stop after selecting N - 1 edges and print their endpoints.
\end{tightitemize}
\end{frame}

\begin{frame}{Walkthrough of the Key State}
\begin{tightitemize}
  \item The solution stores edges and pops them from smallest to largest.
  \item Union-find prevents cycles in the selected road set.
  \item The final list has exactly enough edges to connect all villages.
\end{tightitemize}
\end{frame}

\begin{frame}{Why the Algorithm Is Correct}
\begin{tightitemize}
  \item Kruskal's cut property says the lightest safe edge can be added to some MST.
  \item Union-find identifies whether an edge is safe by detecting if it would form a cycle.
  \item After N - 1 safe edges, the selected graph is connected and acyclic, so it is an MST.
\end{tightitemize}
\end{frame}

\begin{frame}{Complexity and Implementation Notes}
\small
\sectiontag{Complexity:} O(E log E) time and O(E + V) memory; for this complete graph E is O(V\textasciicircum{}2).

\vspace{0.3cm}
\begin{tightitemize}
  \item Output uses 1-based village numbers.
  \item The diagonal self-edges are harmless if ignored by union-find, but they are not useful MST edges.
\end{tightitemize}
\end{frame}



% BEGIN GENERATED CODE WALKTHROUGH
\begin{frame}[fragile]{Code: Initialize Union-Find}
\scriptsize
\sectiontag{Source lines:} \texttt{lostmap.py:53--82}

\vspace{0.08cm}
\begin{columns}[T,totalwidth=\textwidth]
\begin{column}{0.58\textwidth}
\begin{lstlisting}[style=cpcode]
class UnionFind:

    def __init__(self, numItems: int) -> None:
        if numItems < 1:
            raise ValueError("numItems needs to be greater or equal to 1")

        self._numSet = numItems
        self._parents = list(range(numItems))
        self._ranks = [0] * numItems
        self._sizes = [1] * numItems
\end{lstlisting}
\end{column}
\begin{column}{0.38\textwidth}
\sectiontag{What this chunk does}
\begin{tightitemize}
  \item The MST builder needs a DSU to avoid cycles.
  \item Every village begins in its own component.
  \item Rank and size support efficient merging.
\end{tightitemize}
\end{column}
\end{columns}
\end{frame}

\begin{frame}[fragile]{Code: Find Representatives}
\scriptsize
\sectiontag{Source lines:} \texttt{lostmap.py:84--115}

\vspace{0.08cm}
\begin{columns}[T,totalwidth=\textwidth]
\begin{column}{0.58\textwidth}
\begin{lstlisting}[style=cpcode]
def findSet(self, element: int) -> int:
    children = []
    toFind = element

    while toFind != self._parents[toFind]:
        children.append(toFind)
        toFind = self._parents[toFind]

    for c in children:
        self._parents[c] = toFind

    return toFind
\end{lstlisting}
\end{column}
\begin{column}{0.38\textwidth}
\sectiontag{What this chunk does}
\begin{tightitemize}
  \item findSet walks to the representative of a component.
  \item The visited nodes are compressed to point directly to the representative.
  \item This keeps repeated Kruskal checks fast.
\end{tightitemize}
\end{column}
\end{columns}
\end{frame}

\begin{frame}[fragile]{Code: Merge Components}
\scriptsize
\sectiontag{Source lines:} \texttt{lostmap.py:117--155}

\vspace{0.08cm}
\begin{columns}[T,totalwidth=\textwidth]
\begin{column}{0.58\textwidth}
\begin{lstlisting}[style=cpcode]
def isSameSet(self, element1: int, element2: int) -> bool:
    return self.findSet(element1) == self.findSet(element2)

def unionSet(self, element1: int, element2: int) -> None:
    s1 = self.findSet(element1)
    s2 = self.findSet(element2)

    if s1 == s2:
        return None

    if self._ranks[s1] > self._ranks[s2]:
        s1, s2 = s2, s1
    elif self._ranks[s1] == self._ranks[s2]:
        self._ranks[s2] += 1

    self._parents[s1] = s2

    self._numSet -= 1
    self._sizes[s2] += self._sizes[s1]
\end{lstlisting}
\end{column}
\begin{column}{0.38\textwidth}
\sectiontag{What this chunk does}
\begin{tightitemize}
  \item isSameSet checks whether an edge would create a cycle.
  \item unionSet links two different components by rank.
  \item Component sizes and the component count are updated after a successful merge.
\end{tightitemize}
\end{column}
\end{columns}
\end{frame}

\begin{frame}[fragile]{Code: Run Kruskal}
\scriptsize
\sectiontag{Source lines:} \texttt{lostmap.py:177--205}

\vspace{0.08cm}
\begin{columns}[T,totalwidth=\textwidth]
\begin{column}{0.58\textwidth}
\begin{lstlisting}[style=cpcode]
def kruskal(
    numVertex: int, edgeList: list[tuple[int, int, int]]
) -> tuple[list[tuple[int, int, int]], int]:
    edges = sorted(edgeList, key=lambda x: -x[0])
    ufds = UnionFind(numVertex)
    edgeListOutput = []
    totalMSTweight = 0
    while edges and len(edgeListOutput) < numVertex - 1:
        w, u, v = edges.pop()
        ru = ufds.findSet(u)
        rv = ufds.findSet(v)
        if ru != rv:
            ufds.unionSet(ru, rv)
            edgeListOutput.append((w, u, v))
            totalMSTweight += w
    return (edgeListOutput, totalMSTweight)
\end{lstlisting}
\end{column}
\begin{column}{0.38\textwidth}
\sectiontag{What this chunk does}
\begin{tightitemize}
  \item All candidate roads are sorted by weight.
  \item The smallest remaining edge is accepted only if it joins different components.
  \item The loop stops after n - 1 edges, exactly the size of a spanning tree.
\end{tightitemize}
\end{column}
\end{columns}
\end{frame}

\begin{frame}[fragile]{Code: Build Edges from the Table}
\scriptsize
\sectiontag{Source lines:} \texttt{lostmap.py:208--218}

\vspace{0.08cm}
\begin{columns}[T,totalwidth=\textwidth]
\begin{column}{0.58\textwidth}
\begin{lstlisting}[style=cpcode]
lines = sys.stdin.read().strip().splitlines()
grid = [list(map(int, line.split())) for line in lines[1:]]
numVertex = int(lines[0])
output = []
edgeList = []
for row in range(numVertex):
    for col in range(row, numVertex):
        edgeList.append((grid[row][col], row, col))
for edge in kruskal(numVertex, edgeList)[0]:
    output.append(f"{edge[1]+1} {edge[2]+1}\n")
sys.stdout.write("".join(output))
\end{lstlisting}
\end{column}
\begin{column}{0.38\textwidth}
\sectiontag{What this chunk does}
\begin{tightitemize}
  \item The matrix is converted into weighted edges between villages.
  \item Only row <= col is used to avoid duplicating undirected edges.
  \item Kattis expects village numbers starting from 1, so output adds 1.
\end{tightitemize}
\end{column}
\end{columns}
\end{frame}
% END GENERATED CODE WALKTHROUGH

\begin{frame}{Source}
\small
\sectiontag{Solution file:} \texttt{practice\_contests/day\_4/lostmap.py}

\vspace{0.3cm}
\sectiontag{Problem:} \url{https://open.kattis.com/problems/lostmap}

\vspace{0.45cm}
Use this deck with the source code open beside it. The slides explain the reasoning; the code shows the exact input handling and output formatting.
\end{frame}

\end{document}
